\section{Differential variations and the quantitative planar limit}
\label{sec:differential-limit}

\subsection{Constructive differentiable families}

A differentiable real name in this section means a fixed primitive-recursive
procedure for rational approximations to \(f(t)\) and \(f'(t)\), together with
a primitive-recursive modulus \(\delta(n)\) such that every rational
\(0<\abs h\le2^{-\delta(n)}\) satisfies
\begin{equation}
 \abs{\frac{f(t+h)-f(t)}h-f'(t)}\le2^{-n}.                 \tag{9.1}
\end{equation}
The product and chain rules are finite error estimates on bounded intervals.
The ProveIt trigonometric packet proves the derivative names
\((\sin x)'=\cos x\) and \((\cos x)'=-\sin x\) from the finite exponential
series, with explicit tail schedules.  Thus differentiating the cosine laws is
legitimate in \(\ipra\) for every fixed primitive-recursive family.

\begin{theorem}[Differential variation formulas]
\label{thm:differential}
Along a differentiable family of proper spherical triangles,
\begin{align}
 da&=\cos C\,db+\cos B\,dc+\sin b\sin C\,dA,              \tag{9.2}\\
 db&=\cos A\,dc+\cos C\,da+\sin c\sin A\,dB,              \tag{9.3}\\
 dc&=\cos B\,da+\cos A\,db+\sin a\sin B\,dC.              \tag{9.4}
\end{align}
Here each symbol denotes the derivative name at the fixed parameter value.
\end{theorem}

\begin{proof}
Differentiate the cosine law for \(a\):
\begin{align*}
 -\sin a\,da
 &=(-\sin b\cos c+\cos b\sin c\cos A)\,db\\
 &\quad+(-\cos b\sin c+\sin b\cos c\cos A)\,dc
       -\sin b\sin c\sin A\,dA.
\end{align*}
The cyclic five-part identities give
\begin{align*}
 \sin a\cos C&=\sin b\cos c-\cos b\sin c\cos A,\\
 \sin a\cos B&=\cos b\sin c-\sin b\cos c\cos A.
\end{align*}
The sine law gives
\(
 \sin b\sin c\sin A=\sin a\sin b\sin C
\).
Substitution yields
\[
 \sin a\,da=
 \sin a\bigl(\cos C\,db+\cos B\,dc+\sin b\sin C\,dA\bigr).
\]
Since \(\sin a\) has a positive apartness witness, cancellation proves (9.2).
The other formulas are cyclic.  At the finite-precision level the proof is a
sum of product-rule estimates followed by the same polynomial simplification;
no infinitesimal objects are added to the language of arithmetic.
\end{proof}

The equations may be solved for any three unknown differentials once the other
three are supplied and the relevant determinant is positively separated from
zero.  Constructively, matrix inversion is therefore guarded by an explicit
rational lower bound on that determinant.

\subsection{A uniform small-triangle estimate}

The usual statement that the spherical cosine law tends to the planar cosine
law can be made quantitative.

\begin{theorem}[Planar cosine law with an explicit remainder]
\label{thm:planar-limit}
Suppose
\(0\le a,b,c\le\varepsilon\le1\) and the exact spherical cosine law holds.
Then
\begin{equation}
 \abs{c^2-a^2-b^2+2ab\cos C}\le2\varepsilon^4.             \tag{9.5}
\end{equation}
Cyclic variants hold for the other sides.
\end{theorem}

\begin{proof}
The finite Taylor bounds, valid constructively on \([0,1]\), are
\begin{align}
 \abs{t^2-2(1-\cos t)}&\le\frac{t^4}{12},                 \tag{9.6}\\
 0\le t-\sin t&\le\frac{t^3}{6}.                          \tag{9.7}
\end{align}
Put \(F(t)=2(1-\cos t)\).  The spherical cosine law implies
\begin{equation}
 F(c)=F(a)+F(b)-2(1-\cos a)(1-\cos b)-2\sin a\sin b\cos C.
                                                                    \tag{9.8}
\end{equation}
Therefore
\begin{align*}
 &c^2-a^2-b^2+2ab\cos C\\
 &=[c^2-F(c)]-[a^2-F(a)]-[b^2-F(b)]\\
 &\quad-2(1-\cos a)(1-\cos b)
       +2(ab-\sin a\sin b)\cos C.
\end{align*}
The first line of errors has absolute value at most
\(3\varepsilon^4/12\).  Since \(1-\cos t\le t^2/2\), the next term is at most
\(\varepsilon^4/2\).  Finally,
\[
 \abs{ab-\sin a\sin b}
 \le a\abs{b-\sin b}+\sin b\abs{a-\sin a}
 \le\frac{\varepsilon^4}{3},
\]
so the last term is at most \(2\varepsilon^4/3\).  Their sum is
\(17\varepsilon^4/12<2\varepsilon^4\), proving (9.5).
All Taylor estimates come from finite alternating sums with primitive-recursive
remainder indices.
\end{proof}

For a sphere of radius \(R\), let \(\ell_a=Ra\), etc., and assume
\(\ell_a,\ell_b,\ell_c\le L\le R\).  Multiplying (9.5) by \(R^2\) gives
\begin{equation}
 \abs{\ell_c^2-\ell_a^2-\ell_b^2+2\ell_a\ell_b\cos C}
 \le\frac{2L^4}{R^2}.                                      \tag{9.9}
\end{equation}
Thus the curvature correction is explicitly quadratic in \(1/R\).

\begin{corollary}[Planar sine-law cross relation]
Under the same small-side hypothesis,
\begin{equation}
 \abs{a\sin B-b\sin A}\le\frac{\varepsilon^3}{3},          \tag{9.10}
\end{equation}
with cyclic companions.
\end{corollary}

\begin{proof}
The exact spherical sine law says
\(\sin a\sin B=\sin b\sin A\).  Hence
\[
 a\sin B-b\sin A
 =(a-\sin a)\sin B-(b-\sin b)\sin A,
\]
and (9.7) gives the bound.
\end{proof}
